\documentclass[14 pt]{extarticle}

	\usepackage[frenchb]{babel}
	\usepackage[utf8]{inputenc}  
	\usepackage[T1]{fontenc}
	\usepackage{amssymb}
	\usepackage[mathscr]{euscript}
	\usepackage{stmaryrd}
	\usepackage{amsmath}
\usepackage{multicol}
	\usepackage{tikz}
	\usepackage[all,cmtip]{xy}
	\usepackage{amsthm}
	\usepackage{varioref}
	\usepackage{geometry}
	\geometry{a4paper}
	\usepackage{lmodern}
	\usepackage{hyperref}
	\usepackage{array}
	 \usepackage{fancyhdr}
\renewcommand{\theenumi}{\alph{enumi})}
	\pagestyle{fancy}
	\theoremstyle{plain}
	\fancyfoot[C]{} 
	\fancyhead[L]{Interrogation}
	\fancyhead[R]{16 septembre 2025}\geometry{
 a4paper,
 total={170mm,257mm},
 left=15mm,right=10 mm,
 top=20mm,bottom = 10 mm
 }
	
	
	\title{Interrogation Chapitre 1}
	\date{}
	\begin{document}

\begin{center}{\Large Chapitre 1 - Rappels de  numération}\\ 
 \end{center}
  
  \subsection*{Exercice 1 (6 points)}
  
  On considère le nombre $908\ 723, 147$
  
  \begin{enumerate}
  \item Donner son chiffre des dizaines. 
  \item Donner son chiffre des millièmes. 
  \item Donner son chiffre des milliers. 
  \item Recopier la phrase suivante : « Le chiffre des \ldots est égal au chiffre des \ldots »
  \item Écrire ce nombre comme une fraction décimale. 
  \item Écrire ce nombre comme un nombre mixte.
  \end{enumerate}
  
  \subsection*{Exercice 2 (5 points)}
  Donner l'écriture décimale des nombres suivants. 
  \begin{enumerate}
  \begin{multicols}{2}
  \item $\displaystyle\frac{539}{10}$
  \item $\displaystyle\frac{124}{10~000}$
  \item $\displaystyle432+ \frac{95}{100}$
  \item $\displaystyle64+\frac{543}{100}$
  \item $\displaystyle1\times 100 + 6 \times 1 + \frac5{10} + \frac3{100} + \frac4{10~000}$
  \end{multicols}
  \end{enumerate}
  
  \subsection*{Exercice 3 (5 points)}
  Donner une fraction décimale égale à : 
  \begin{enumerate}
 \begin{multicols}{2} \item $\displaystyle94,134$
  \item $\displaystyle34+ \frac{85}{100}$
  \item $\displaystyle\frac{7}{10} + \frac4{100} + \frac{9}{1000}$
  \item $\displaystyle6 + \frac{16}{10}$. 
  \item $\displaystyle\frac{12}{10}+\frac{63}{100}+ \frac{93}{1~000}$
  \end{multicols}
  \end{enumerate}
  
  \subsection*{Exercice 4 (4 points)}
Écrire sous la forme d'un nombre mixte. 
\begin{enumerate}\begin{multicols}{2}
\item $\displaystyle344,923$
\item $\displaystyle\frac{854}{10}$
\item $\displaystyle\frac{21}{10}+\frac{21}{100}$
\item $\displaystyle843 + \frac{912}{100}$
\end{multicols}
\end{enumerate}

\newpage
\begin{center}{\Large Chapitre 1 - Rappels de  numération}\\ 
 \end{center}
  
  \subsection*{Exercice 1 (6 points)}
  
  On considère le nombre $908\ 723, 184$
  
  \begin{enumerate}
  \item Donner son chiffre des dizaines. 
  \item Donner son chiffre des millièmes. 
  \item Donner son chiffre des milliers. 
  \item Recopier la phrase suivante : « Le chiffre des \ldots est égal au chiffre des \ldots »
  \item Écrire ce nombre comme une fraction décimale. 
  \item Écrire ce nombre comme un nombre mixte.
  \end{enumerate}
  
  \subsection*{Exercice 2 (5 points)}
  Donner l'écriture décimale des nombres suivants. 
  \begin{enumerate} \begin{multicols}{2}
  \item $\displaystyle\frac{519}{10}$
  \item $\displaystyle\frac{123}{10~000}$
  \item $\displaystyle432+ \frac{95}{100}$
  \item $\displaystyle64+\frac{543}{100}$
  \item $\displaystyle1\times 100 + 6 \times 1 + \frac5{10} + \frac3{100} + \frac4{10~000}$
  \end{multicols}
  \end{enumerate}
  
  \subsection*{Exercice 3 (5 points)}
  Donner une fraction décimale égale à : 
  \begin{enumerate} \begin{multicols}{2}
  \item $\displaystyle46,134$
  \item $\displaystyle34+ \frac{85}{100}$
  \item $\displaystyle\frac{7}{10} + \frac4{100} + \frac{9}{1000}$
  \item $\displaystyle4 + \frac{16}{10}$. 
  \item $\displaystyle\frac{12}{10}+\frac{63}{100}+ \frac{93}{1~000}$
  \end{multicols}
  \end{enumerate}
  
  \subsection*{Exercice 4 (4 points)}
Écrire sous la forme d'un nombre mixte. 
\begin{enumerate} \begin{multicols}{2}
\item $\displaystyle843,91$
\item $\displaystyle\frac{942}{10}$
\item $\displaystyle\frac{91}{10}+\frac{21}{100}$
\item $\displaystyle123 + \frac{912}{100}$
\end{multicols}
\end{enumerate}

 	\end{document}
